\subsection{限制对偶上的余代数结构}

对任意的 \(a\)（\(\mathbf{A}\) 中的），我们用 \(\eta(a)\) 表示线性型 \(f \mapsto f(a)\)（\(\Theta(\mathbf{A})\) 上的）。这样我们就定义了一个 K-线性映射 \(\eta\)，从 \(\mathbf{A}\) 到对偶 \(\Theta(\mathbf{A})^*\)（向量空间 \(\Theta(\mathbf{A})\) 的对偶）。令 \(\varepsilon = \eta(1)\)。

命题 8. —— 在向量空间 \(\Theta(A)\) 上存在唯一的余代数结构（III, §11, No.~1, 第 574 页），使得映射 \(\eta: A \to \Theta(A)^*\) 是从 \(A\) 到 \(\Theta(A)\) 的对偶代数（III, §11, No.~2, 第 579 页）的同态。余代数 \(\Theta(A)\) 是余结合的，且以 \(\varepsilon\) 为余单位。

对每个整数 \(n \geqslant 1\)，考虑 K-线性映射 \(j_{n}\)（从 \(\Theta(\mathrm{A})^{\otimes n}\) 到 \(A^{\otimes n}\) 的对偶），它由公式

\[
\langle j _ {n} \left(f _ {1} \otimes \dots \otimes f _ {n}\right), a _ {1} \otimes \dots \otimes a _ {n} \rangle = \prod_ {i = 1} ^ {n} \left\langle f _ {i}, a _ {i} \right\rangle\tag{13}
\]

刻画，其中 \((a_{i})\) 属于 \(\mathrm{A}^n\)，\((f_i)\) 属于 \(\Theta (\mathrm{A})^n\)。由 II, §7, No.~7, 第 308 页的命题 16 (ii)，映射 \(j_{n}\) 是单射。用 \(m_{\mathrm{K}}:\mathrm{K}\otimes \mathrm{K}\to \mathrm{K}\) 与 \(m_{\mathrm{A}}:\mathrm{A}\otimes \mathrm{A}\rightarrow \mathrm{A}\) 分别表示由 K 与 A 中的乘法导出的映射。对 \(f,g\in \Theta (\mathrm{A})\) 与 \(a,b\in \mathrm{A}\)，我们有

\[
\langle j _ {2} (f \otimes g), a \otimes b \rangle = m _ {\mathrm{K}} \circ (\eta (a) \otimes \eta (b)) (f \otimes g).
\]

因此我们有

\[
\langle j _ {2} (t), a \otimes b \rangle = m _ {\mathrm{K}} \circ (\eta (a) \otimes \eta (b)) (t)\tag{14}
\]

对每个 \(t\in \Theta (\mathrm{A})\otimes \Theta (\mathrm{A})\) 成立。

引理 2. —— 设 \(c: \Theta(A) \to \Theta(A) \otimes \Theta(A)\) 是一个 K-线性映射。那么 \(\eta\) 是从 A 到余代数 \((\Theta(A), c)\) 的对偶代数的同态，当且仅当下面的图交换：

\[
\begin{array}{c} \Theta (\mathrm{A}) \xrightarrow {c} \Theta (\mathrm{A}) \otimes \Theta (\mathrm{A}) \\ \Biggl \downarrow_ {j _ {1}} \qquad \qquad \qquad \qquad \Biggl \downarrow_ {j _ {2}} \\ \mathrm{A} ^ {*} \xrightarrow {^ t m _ {\mathrm{A}}} (\mathrm{A} \otimes \mathrm{A}) ^ {*}. \end{array}\tag{15}
\]

事实上，\(\eta\) 是从 A 到余代数 \((\Theta (\mathrm{A}),c)\) 的对偶代数的同态，当且仅当 \(\eta (ab) = m_{\mathrm{K}}\circ (\eta (a)\otimes \eta (b))\circ c\) 对一切 \(a,b\in \mathrm{A}\) 成立，也就是

\[
\eta (a b) (f) = m _ {\mathrm{K}} \circ (\eta (a) \otimes \eta (b)) (c (f))
\]

对 \(a, b \in \mathrm{A}\) 与 \(f \in \Theta(\mathrm{A})\) 成立。而今，我们有

\[
\eta (a b) (f) = f (a b) = \left\langle {} ^ {t} m _ {\mathrm{A}} \left(j _ {1} (f)\right), a \otimes b \right\rangle
\]

而由 (14)，

\[
m _ {K} \circ (\eta (a) \otimes \eta (b)) (c (f)) = \langle j _ {2} (c (f)), a \otimes b \rangle
\]

对一切 \(a, b \in \mathrm{A}\) 与每个 \(f \in \Theta(\mathrm{A})\) 成立。引理得证。

由于 \(j_{2}\) 是单射，至多存在一个使上图交换的线性映射 c。为证明其存在性，必须证明 \(^{t}m \circ j_{1}\) 的像包含于 \(j_{2}\) 的像。换言之，必须证明：对 \(\Theta(\mathrm{A})\) 的每个元素 f，存在一个自然数 n 以及元素 \(f_{1}^{\prime}, \ldots, f_{n}^{\prime}, f_{1}^{\prime\prime}, \ldots, f_{n}^{\prime\prime}\)（\(\Theta(\mathrm{A})\) 中的元素），满足关系式

\[
f (a b) = \sum_ {i = 1} ^ {n} f _ {i} ^ {\prime} (a) f _ {i} ^ {\prime \prime} (b)\tag{16}
\]

其中 \(a, b \in \mathbf{A}\)。此时我们将有

\[
c (f) = \sum_ {i = 1} ^ {n} f _ {i} ^ {\prime} \otimes f _ {i} ^ {\prime \prime}.\tag{17}
\]

由 VIII, 第 379 页的推论，存在一个以 f 为系数的 K 上有限维左 A-模 E。设 \((e_{1},\ldots,e_{n})\) 是 E 的一组基，\((e_{1}^{*},\ldots,e_{n}^{*})\) 是对偶基，x 是 E 的一个元素，\(x^{*}\) 是 \(E^{*}\) 的一个元素且满足 \(f=c_{\mathrm{E}}(x,x^{*})\)。令 \(f_{i}^{\prime}=c_{\mathrm{E}}(e_{i},x^{*})\) 与 \(f_{i}^{\prime\prime}=c_{\mathrm{E}}(x,e_{i}^{*})\)（对 \(i\in[1,n]\)）；对

A 中的 \(a, b\)，我们有

\[
\begin{array}{r l} & {f (a b) = \langle x ^ {*}, a b x \rangle = \langle x ^ {*} a, b x \rangle = \left\langle \sum_ {i} \langle x ^ {*} a, e _ {i} \rangle e _ {i} ^ {*}, b x \right\rangle} \\ & {\qquad = \sum_ {i} \langle x ^ {*} a, e _ {i} \rangle \langle e _ {i} ^ {*}, b x \rangle = \sum_ {i} \langle x ^ {*}, a e _ {i} \rangle \langle e _ {i} ^ {*}, b x \rangle = \sum_ {i} f _ {i} ^ {\prime} (a) f _ {i} ^ {\prime \prime} (b),} \end{array}
\]

从而得到 (16)。

我们来证明 \(c\) 的余结合性。为此，考虑 K-线性映射

\[
c ^ {\prime} = \left(c \otimes 1 _ {\Theta (\mathrm{A})}\right) \circ c \quad \text { and } \quad c ^ {\prime \prime} = \left(1 _ {\Theta (\mathrm{A})} \otimes c\right) \circ c
\]

它们都是从 \(\Theta (\mathrm{A})\) 到 \(\Theta (\mathrm{A})^{\otimes 3}\) 的映射。我们有关系式

\[
\begin{array}{r l} \langle j _ {3} (f \otimes c (g)), a \otimes b \otimes c \rangle & = \langle f, a \rangle \langle j _ {2} \circ c (g), b \otimes c \rangle \\ & = \langle f, a \rangle \langle g, b c \rangle \\ & = \langle j _ {2} (f \otimes g), a \otimes b c \rangle \\ & = \langle^ {t} (\mathrm{Id} _ {\mathrm{A}} \otimes m _ {\mathrm{A}}) \circ j _ {2} (f \otimes g), a \otimes b \otimes c \rangle \end{array}
\]

其中 \(f, g \in \Theta(\mathbf{A})\)，\(a, b, c \in \mathbf{A}\)。由此我们导出下面的图是交换的：

\[
\begin{array}{c} \Theta (\mathrm{A}) \otimes \Theta (\mathrm{A}) \xrightarrow {\operatorname{Id} _ {\Theta (\mathrm{A})} \otimes c} \Theta (\mathrm{A}) \otimes \Theta (\mathrm{A}) \otimes \Theta (\mathrm{A}) \\ \Biggl \downarrow_ {j _ {2}} \\ (\mathrm{A} \otimes \mathrm{A}) ^ {*} \xrightarrow {^ t (\operatorname{Id} _ {\mathrm{A}} \otimes m _ {\mathrm{A}})} (\mathrm{A} \otimes \mathrm{A} \otimes \mathrm{A}) ^ {*}. \end{array}\tag{18}
\]

由于这个图以及 (15) 的交换性，对 \(f \in \Theta(\mathrm{A})\) 与 \(a, a', a'' \in A\)，我们有

\[
\langle j _ {3} \circ c ^ {\prime} (f), a \otimes a ^ {\prime} \otimes a ^ {\prime \prime} \rangle = \langle f, (a a ^ {\prime}) a ^ {\prime \prime} \rangle ;
\]

同理可证关系式

\[
\langle j _ {3} \circ c ^ {\prime \prime} (f), a \otimes a ^ {\prime} \otimes a ^ {\prime \prime} \rangle = \langle f, a (a ^ {\prime} a ^ {\prime \prime}) \rangle
\]

由于 A 中的乘法是结合的，我们有 \(j_{3} \circ c' = j_{3} \circ c''\)，因此 \(c' = c''\)，因为 \(j_{3}\) 是单射。

最后，公式 (16) 与 (17) 蕴含 \(\Theta(\mathrm{A})\) 以 \(\varepsilon\) 为余单位。

注 1. —— 设 \((\mathrm{V},\pi)\) 是代数 A 的一个有限维线性表示。取 V 的一组基 \((e_{1},\ldots,e_{n})\) 以及对偶基 \((e_{1}^{*},\ldots,e_{n}^{*})\)（\(V^{*}\) 的对偶基）。由引理 2 的证明，我们有关系式

\[
c (c _ {\pi} (x, x ^ {*})) = \sum_ {k = 1} ^ {n} c _ {\pi} (e _ {k}, x ^ {*}) \otimes c _ {\pi} (x, e _ {k} ^ {*})\tag{19}
\]

其中 \(x \in V\)，\(x^{*} \in V^{*}\)。对 \(1 \leqslant i, j \leqslant n\)，令 \(\pi_{ij} = c_{\pi}(e_j, e_i^*)\)。对每个 \(a \in A\)，\(\pi(a)\) 关于 V 的基 \((e_1, \ldots, e_n)\) 的矩阵等于 \((\pi_{ij}(a))\)。于是，对 \(1 \leqslant i \leqslant n\) 与 \(1 \leqslant j \leqslant n\)，我们有

\[
c (\pi_ {i j}) = \sum_ {k = 1} ^ {n} \pi_ {i k} \otimes \pi_ {k j}.\tag{20}
\]

定义 3. —— 设 C 是体 K 上的一个余代数，c 是它的余积。C 的余子代数是指 C 的这样的线性子空间 \(C_{1}\)：\(c(\mathrm{C}_{1})\) 包含于 \(C_{1} \otimes_{K} C_{1}\) 在 \(C \otimes_{K} C\) 中的典范像。

设 \(j\) 是 \(C_1\) 到 \(C\) 中的典范单射。按此定义，存在唯一的线性映射 \(c_1: C_1 \to C_1 \otimes_K C_1\)，使得我们有

\[
c \circ j = (j \otimes j) \circ c _ {1};\tag{21}
\]

这一关系式意味着 \(j\) 是从 \((\mathrm{C}_1, c_1)\) 到 \((\mathrm{C}, c)\) 的余代数态射（III, §11, No.~1, 第 574 页）。

若 C 是余结合的，则 \(C_{1}\) 也是余结合的。若 C 是余交换的，则 \(C_{1}\) 也是余交换的。若 C 有余单位 \(\varepsilon\)，则 \(\varepsilon\) 在 \(C_{1}\) 上的限制是 \(C_{1}\) 的一个余单位。

命题 9. —— 设 \(\Theta\) 是 \(\Theta(A)\) 的一个线性子空间。下列性质等价：

\begin{enumerate}
\def\labelenumi{(\roman{enumi})}
\item
  \(\Theta\) 是 \(\Theta(A)\) 的一个 (A, A)-子双模。
\item
  \(\Theta\) 是 \(\Theta(A)\) 的一个余子代数。
\item
  存在 A-模的类的一个遗传集 \(\mathcal{C}\)（这些 A-模都在 \(K\) 上有限维），使得 \(\Theta = \Theta_{\mathcal{C}}(A)\)。
\end{enumerate}

当这些性质成立时，(iii) 中所述的集合 C 是唯一确定的。

最后一个断言由 VIII, 第 388 页的推论得出：集合 C 由满足 \(\Theta_{\mathrm{E}}(\mathrm{A})\) 包含于 \(\Theta\) 的 K 上有限维 A-模 E 的类组成。

\begin{enumerate}
\def\labelenumi{(\roman{enumi})}
\setcounter{enumi}{2}
\item
  \(\Rightarrow\) (ii)：设 C 是 K 上有限维 A-模的类的一个遗传集。那么 \(\Theta_{\mathcal{C}}(\mathrm{A})\) 是有向族 \((\Theta_{\mathrm{E}}(\mathrm{A}))_{\mathrm{E}\in\mathcal{C}}\) 的并。由于 \(\Theta_{\mathrm{E}}(\mathrm{A})\) 是 \(\Theta(\mathrm{A})\) 的余子代数（对每个 \(E\in C\)）（VIII, 第 390 页，公式 (19)），故对 \(\Theta_{\mathcal{C}}(\mathrm{A})\) 亦然。
\item
  \(\Rightarrow\) (i)：设 \(f\in \Theta (\mathrm{A})\)。设 \(f_{1}^{\prime},\ldots ,f_{n}^{\prime},f_{1}^{\prime \prime},\ldots ,f_{n}^{\prime \prime}\) 是 \(\Theta (\mathrm{A})\) 中满足 \(c(f) = \sum f_i^\prime \otimes f_i''\) 的元素。对 A 中的 \(a,b\)，我们有 \(f(ab) = \sum f_i'(a)f_i''(b)\)，因此
\end{enumerate}

\[
b f = \sum_ {i = 1} ^ {n} f _ {i} ^ {\prime \prime} (b) f _ {i} ^ {\prime}, \quad f a = \sum_ {i = 1} ^ {n} f _ {i} ^ {\prime} (a) f _ {i} ^ {\prime \prime}
\]

（VIII, 第 375 页，公式 (3) 与 (4)）。因此，\(\Theta(A)\) 的余子代数是一个 \((A, A)\)-子双模。

\begin{enumerate}
\def\labelenumi{(\roman{enumi})}
\tightlist
\item
  \(\Rightarrow\) (iii)：设 \(\Theta\) 是 \(\Theta(\mathrm{A})\) 的一个 (A, A)-子双模；令 C 是满足 \(\Theta_{\mathrm{E}}(\mathrm{A})\) 包含于 \(\Theta\) 的 K 上有限维 A-模 E 的类所成的集合。集合 C 是遗传的（VIII, 第 377 页，注 2），且按构造有 \(\Theta_{\mathcal{C}}(\mathrm{A}) \subset \Theta\)。设 \(f \in \Theta\) 且 E = Af。那么 E 在 K 上有限维。于是，E 上的每个线性型都形如 \(u_{a}: g \mapsto g(a)\)，其中 a 属于 A（II, §7, No.~5, 第 302 页，推论 2）。而今，对 \(a, b, x \in A\)，我们有
\end{enumerate}

\[
c _ {\mathrm{E}} (a f, u _ {b}) (x) = \langle u _ {b}, x a f \rangle = f (b x a) = a f b (x),
\]

从而 \(c_{\mathrm{E}}(af, u_{b}) = afb\)。于是我们有 \(\Theta_{\mathrm{E}}(\mathrm{A}) \subset \Theta\)。因此 A-模 E 是 C 型的，且 f 是它的一个系数。于是我们有 \(\Theta \subset \Theta_{\mathcal{C}}(\mathrm{A})\)，最终得 \(\Theta = \Theta_{\mathcal{C}}(\mathrm{A})\)。

注 2. —— 设 \(\Theta\) 是 \(\Theta(A)\) 的一个余子代数。我们赋予 \(K\) 离散拓扑，并赋予代数 \(A\) 使映射 \(f: A \to K\)（其中 \(f\) 取遍 \(\Theta\)）都连续的最粗拓扑（Gen. Top., I, §2, No.~2, 第 175 页）。这个拓扑赋予 \(A\) 一个拓扑 \(K\)-模的结构。双边理想 \(\mathfrak{a}\)（\(A\) 的双边理想）是开的，当且仅当它具有有限余维数且其正交 \(\mathfrak{a}'\) 包含于 \(\Theta\)。由 VIII, 第 379 页的命题 3，\(A\) 的开双边理想构成 \(0\) 在 \(A\) 中的一个基本邻域系。因此 \(A\) 上的拓扑与其环结构相容。

设 C 是满足 \(\Theta = \Theta_{C}\) 的 K 上有限维 A-模的类的遗传集（VIII, 第 391 页，命题 9）。设 E 是 K 上有限维的左 A-模。我们赋予它离散拓扑。下列性质等价：

\begin{enumerate}
\def\labelenumi{(\roman{enumi})}
\item
  A-模 E 是 C 型的。
\item
  A-模 E 的零化理想在 A 中是开的。
\item
  从 A × E 到 E 的映射 \((a, x) \mapsto ax\) 是连续的。
\end{enumerate}

最后一个性质意味着 \(\mathrm{E}\) 是一个拓扑 A-模。

设 \(\Theta^{*}\) 是余代数 \(\Theta\) 的对偶代数。我们赋予 \(\Theta^{*}\) 使映射 \(\varphi \mapsto \varphi(u)\)（从 \(\Theta^{*}\) 到 K，其中 \(u\) 取遍 \(\Theta\)）都连续的最粗拓扑。代数 \(\Theta^{*}\) 上的拓扑与 \(\Theta^{*}\) 上的加群结构相容。\(\Theta^{*}\) 中形如 \(\Theta_{\mathrm{E}}(\mathrm{A})\)（其中 E 为 \(\mathcal{C}\) 型 A-模）的集合的正交构成 0 的一个基本邻域系。而今，这样的集合是 \(\Theta\) 的余子代数，故其正交是 \(\Theta^{*}\) 的理想。因此 \(\Theta^{*}\) 上的拓扑与其环结构相容。典范代数同态 \(\eta: A \to \Theta^{*}\)（它把 \(a \in A\) 映为线性型 \(f \mapsto f(a)\)，即 \(\Theta\) 上的线性型）定义了从 A 的完备豪斯多夫空间 \(\widehat{A}\)（Gen. Top., II, §3, No.~7, 第 191 页）到 \(\Theta^{*}\) 的一个同构。

